% BHU PYQ -- Transcribed & Corrected
% Paper: BHUVA-22 / BHUVA22: Environmental Science (End Term)
% Exam: Under Graduate Semester II Examination, 2025-26 (NEP)
% Category: Value Added Course (VAC)
% Department: Environmental Science

\documentclass[11pt,a4paper]{article}
\usepackage[a4paper,top=2.5cm,bottom=2.5cm,left=2.8cm,right=2.8cm,headheight=14pt]{geometry}
\usepackage{amsmath,amssymb}
\usepackage[T1]{fontenc}
\usepackage[utf8]{inputenc}
\usepackage{lmodern,microtype,fancyhdr,enumitem,hyperref,lastpage,parskip}

\hypersetup{
    colorlinks=true,
    urlcolor=black,
    linkcolor=black,
    pdftitle={BHUVA22: Environmental Science -- BHU Sem II 2025-26 (NEP VAC)}
}

\pagestyle{fancy}
\fancyhf{}
\renewcommand{\headrulewidth}{0.4pt}
\renewcommand{\footrulewidth}{0.4pt}
\fancyhead[L]{\small   \textit{Banaras Hindu University}}
\fancyhead[R]{\small   \textit{BHUVA-22: Environmental Science (VAC)}}
\fancyfoot[C]{\small Page  \thepage\ of \pageref{LastPage}}


\newlist{parts}{enumerate}{2}
\setlist[parts,1]{label=(\arabic*),leftmargin=*,itemsep=3pt,topsep=2pt}

\newcommand{\pts}[1]{\hfill   \textit{[#1]}}


\begin{document}


\begin{center}
  {\large\bfseries BANARAS HINDU UNIVERSITY}\\[2pt]
  {
\normalsize Under Graduate --- Semester II Examination, 2025-26 (NEP)}\\[6pt]
  \rule{0.8\linewidth}{0.5pt}\\[6pt]
  {\large\bfseries Value Added Course (VAC)}\\[4pt]
  {\large\bfseries Paper Code: BHUVA-22 : Environmental Science (End Term)}\\[4pt]
  {
\normalsize Time: 1.5 Hours\hfill Maximum Marks: 70}\\[2pt]
  {\small REF NO. 260718 / 260719}
\end{center}

\rule{\linewidth}{0.4pt}

\smallskip



\noindent   \textit{Instructions: Answer all 70 Multiple Choice Questions. Each question carries equal marks (1 mark each).}

\bigskip


\begin{enumerate}[label=   \textbf{Question \arabic*.}]

    \item Environmental Science is best defined as:\pts{1}
    
\begin{parts}
        \item The study of human society only
        \item The study of physical and biological components of the environment and their interactions
        \item The study of industrial processes
        \item The study of extraterrestrial bodies
    \end{parts}

    \medskip

    \item Which of the following natural resources is renewable?\pts{1}
    
\begin{parts}
        \item Petroleum
        \item Solar Energy
        \item Coal
        \item Natural Gas
    \end{parts}

    \medskip

    \item The main goal of the Paris Agreement (2015) is to:\pts{1}
    
\begin{parts}
        \item Reduce greenhouse gas emissions to limit global warming below 2$^\circ$C
        \item Protect endangered species from illegal trade
        \item Promote space exploration
        \item Increase coal mining
    \end{parts}

    \medskip

    \item What is the full form of GMO?\pts{1}
    
\begin{parts}
        \item Green Modified Organism
        \item Genetically Modified Organism
        \item Globally Managed Object
        \item Genetically Managed Option
    \end{parts}

    \medskip

    \item Which human activity is the primary cause of deforestation?\pts{1}
    
\begin{parts}
        \item Agricultural expansion and logging
        \item Ecotourism
        \item Afforestation
        \item Rainwater harvesting
    \end{parts}

    \medskip

    \item What does GMO stand for in biotechnology?\pts{1}
    
\begin{parts}
        \item Green Modified Organism
        \item Genetically Modified Organism
        \item Globally Managed Object
        \item Genetically Managed Option
    \end{parts}

    \medskip

    \item The largest tropical rainforest in the world is:\pts{1}
    
\begin{parts}
        \item Amazon Rainforest
        \item Congo Rainforest
        \item Daintree Rainforest
        \item Southeast Asian Rainforest
    \end{parts}

    \medskip

    \item Which energy source was most responsible for early industrialization?\pts{1}
    
\begin{parts}
        \item Solar energy
        \item Wind energy
        \item Coal
        \item Nuclear energy
    \end{parts}

    \medskip

    \item Which organization is responsible for disaster management in India?\pts{1}
    
\begin{parts}
        \item NHAI
        \item NDMA
        \item RBI
        \item ICAR
    \end{parts}

    \medskip

    \item What do environmental laws regulate?\pts{1}
    
\begin{parts}
        \item Air quality
        \item Water quality
        \item Wildlife protection
        \item All of the above
    \end{parts}

    \medskip

    \item Who is considered the founder of modern environmentalism?\pts{1}
    
\begin{parts}
        \item Rachel Carson
        \item Aldo Leopold
        \item John Muir
        \item Henry David Thoreau
    \end{parts}

    \medskip

    \item The Chernobyl nuclear disaster occurred in which year?\pts{1}
    
\begin{parts}
        \item 1979
        \item 1986
        \item 2011
        \item 1999
    \end{parts}

    \medskip

    \item Which significant environmental awareness day is celebrated globally on June 5th?\pts{1}
    
\begin{parts}
        \item International Labour Day
        \item International Day for Biological Diversity
        \item World Environment Day
        \item Earth Day
    \end{parts}

    \medskip

    \item The solid outer layer of the Earth is called:\pts{1}
    
\begin{parts}
        \item Atmosphere
        \item Hydrosphere
        \item Lithosphere
        \item Biosphere
    \end{parts}

    \medskip

    \item Which of these is not included in the Red List of Threatened Species?\pts{1}
    
\begin{parts}
        \item Vulnerable
        \item Near Threatened
        \item Extinct
        \item Less Endangered
    \end{parts}

    \medskip

    \item Which of these is a book related to pesticides and environmental pollution written by Rachel Carson?\pts{1}
    
\begin{parts}
        \item A poison like no other
        \item Nanoworld
        \item Silent Spring
        \item The Substance of Civilization
    \end{parts}

    \medskip

    \item How many Sustainable Development Goals (SDGs) were established by the UN?\pts{1}
    
\begin{parts}
        \item 10
        \item 15
        \item 17
        \item 20
    \end{parts}

    \medskip

    \item Which historical industrial disaster occurred in Bhopal, India in 1984?\pts{1}
    
\begin{parts}
        \item Bhopal Gas Tragedy (Methyl Isocyanate leak)
        \item Chernobyl Nuclear Accident
        \item Exxon Valdez Oil Spill
        \item Fukushima Disaster
    \end{parts}

    \medskip

    \item The term for all life forms and their global habitats on Earth is:\pts{1}
    
\begin{parts}
        \item Biosphere
        \item Lithosphere
        \item Hydrosphere
        \item Atmosphere
    \end{parts}

    \medskip

    \item The term `Ecosystem' was coined by:\pts{1}
    
\begin{parts}
        \item Ernst Haeckel
        \item A. G. Tansley
        \item Eugene Odum
        \item Ramdeo Misra
    \end{parts}

    \medskip

    \item Pandemics are best described as:\pts{1}
    
\begin{parts}
        \item Disease limited to a single village
        \item Disease affecting a city
        \item Disease spreading across countries or continents
        \item Non-infectious disease
    \end{parts}

    \medskip

    \item Which layer of the atmosphere is the coldest?\pts{1}
    
\begin{parts}
        \item Troposphere
        \item Stratosphere
        \item Mesosphere
        \item Thermosphere
    \end{parts}

    \medskip

    \item IUCN stands for:\pts{1}
    
\begin{parts}
        \item International Union for Conservation of Natural Resources
        \item International Unit of Conventional Nomenclature
        \item International Union of Carbohydrates and Nucleic Acids
        \item International Union for Conservation of Nature
    \end{parts}

    \medskip

    \item Noise pollution is measured in:\pts{1}
    
\begin{parts}
        \item Hertz
        \item Decibel (dB)
        \item Joule
        \item Pascal
    \end{parts}

    \medskip

    \item Global warming is mainly caused by an increase in atmospheric concentrations of:\pts{1}
    
\begin{parts}
        \item Oxygen
        \item Nitrogen
        \item Carbon Dioxide ($ \text{CO}_2$)
        \item Ozone
    \end{parts}

    \medskip

    \item The Neolithic period is most closely associated with which major human-environment interaction?\pts{1}
    
\begin{parts}
        \item Nomadic hunting and food gathering
        \item Beginning of agriculture and domestication of animals
        \item Industrial manufacturing
        \item Urbanization and smart cities
    \end{parts}

    \medskip

    \item Which of the following is an abiotic component of an ecosystem?\pts{1}
    
\begin{parts}
        \item Bacteria
        \item Algae
        \item Temperature
        \item Plants
    \end{parts}

    \medskip

    \item Which environmental issue results from excessive burning of fossil fuels?\pts{1}
    
\begin{parts}
        \item Soil erosion
        \item Acid rain
        \item Earthquakes
        \item Tsunami
    \end{parts}

    \medskip

    \item Excessive noise pollution mainly affects:\pts{1}
    
\begin{parts}
        \item Soil fertility
        \item Human health and hearing
        \item Water cycle
        \item Atmospheric pressure
    \end{parts}

    \medskip

    \item Natural capital degradation refers to:\pts{1}
    
\begin{parts}
        \item Soil erosion
        \item Aquifer depletion
        \item Water pollution
        \item All of the above
    \end{parts}

    \medskip

    \item What is the full form of EPA?\pts{1}
    
\begin{parts}
        \item Environmental Provisioning Agency
        \item Environmental Protection Authority
        \item Environmental Prevention Authority
        \item Environmental Protection Agency
    \end{parts}

    \medskip

    \item One major consequence of global climate change is:\pts{1}
    
\begin{parts}
        \item Stable sea level
        \item Melting of glaciers and polar ice caps
        \item Increase in forest cover
        \item Reduction in solar radiation
    \end{parts}

    \medskip

    \item Which International Agreement aims specifically to protect the ozone layer?\pts{1}
    
\begin{parts}
        \item Kyoto Protocol
        \item Paris Agreement
        \item Montreal Protocol
        \item Rio Declaration
    \end{parts}

    \medskip

    \item Primary consumers in a food chain are:\pts{1}
    
\begin{parts}
        \item Carnivores
        \item Herbivores
        \item Decomposers
        \item Omnivores
    \end{parts}

    \medskip

    \item Sustainable Development Goals (SDGs) were adopted by the United Nations in:\pts{1}
    
\begin{parts}
        \item 2000
        \item 2010
        \item 2015
        \item 2020
    \end{parts}

    \medskip

    \item Bishnoi Movement in Rajasthan was aimed at protecting:\pts{1}
    
\begin{parts}
        \item Khejri trees and wildlife
        \item Rivers
        \item Mining sites
        \item Agricultural lands
    \end{parts}

    \medskip

    \item Which of the following is a non-renewable energy resource?\pts{1}
    
\begin{parts}
        \item Coal
        \item Solar energy
        \item Wind energy
        \item Hydro energy
    \end{parts}

    \medskip

    \item Rapid urbanization often results in:\pts{1}
    
\begin{parts}
        \item Increased waste generation and infrastructure strain
        \item Improved biodiversity
        \item Unplanned deforestation reduction
        \item Lower energy consumption
    \end{parts}

    \medskip

    \item Ex-situ conservation methods include:\pts{1}
    
\begin{parts}
        \item Botanical gardens and seed banks
        \item National Parks
        \item Wildlife Sanctuaries
        \item Biosphere Reserves
    \end{parts}

    \medskip

    \item For the conservation of wetlands, which international treaty was signed?\pts{1}
    
\begin{parts}
        \item Stockholm Convention
        \item Ramsar Convention
        \item Bonn Convention
        \item Vienna Convention
    \end{parts}

    \medskip

    \item Urbanization refers to:\pts{1}
    
\begin{parts}
        \item Decline in city population
        \item Growth of rural population
        \item Increase in urban population and expansion of cities
        \item Development of traditional agriculture
    \end{parts}

    \medskip

    \item Decomposers in an ecosystem mainly include:\pts{1}
    
\begin{parts}
        \item Plants
        \item Animals
        \item Bacteria and Fungi
        \item Algae
    \end{parts}

    \medskip

    \item Which biome is known for the highest biodiversity and dense tree canopy?\pts{1}
    
\begin{parts}
        \item Tropical Rainforest
        \item Tundra
        \item Desert
        \item Temperate Grassland
    \end{parts}

    \medskip

    \item Biodiversity represents:\pts{1}
    
\begin{parts}
        \item Variety of ecosystems
        \item Genetic variation within species
        \item Number and variety of species in a habitat
        \item All of the above
    \end{parts}

    \medskip

    \item Which terrestrial ecosystem is characterized by permafrost and absence of trees?\pts{1}
    
\begin{parts}
        \item Desert
        \item Grassland
        \item Tundra
        \item Tropical Forest
    \end{parts}

    \medskip

    \item The boundary between troposphere and stratosphere is called:\pts{1}
    
\begin{parts}
        \item Tropopause
        \item Stratopause
        \item Mesopause
        \item Thermopause
    \end{parts}

    \medskip

    \item High concentration of particulate matter ($ \text{PM}_{2.5}$) in the air primarily causes:\pts{1}
    
\begin{parts}
        \item Respiratory diseases
        \item Water pollution
        \item Soil erosion
        \item Noise pollution
    \end{parts}

    \medskip

    \item Which energy source produces zero carbon emissions during operation?\pts{1}
    
\begin{parts}
        \item Solar and Wind power
        \item Coal power plants
        \item Diesel generators
        \item Natural gas
    \end{parts}

    \medskip

    \item Which acid is primarily responsible for Acid Rain?\pts{1}
    
\begin{parts}
        \item Sulfuric acid ($\ch{H2SO4}$) and Nitric acid ($\ch{HNO3}$)
        \item Hydrochloric acid ($\ch{HCl}$)
        \item Acetic acid ($\ch{CH3COOH}$)
        \item Citric acid
    \end{parts}

    \medskip

    \item The primary goal of waste management hierarchy is:\pts{1}
    
\begin{parts}
        \item Reduce, Reuse, Recycle
        \item Incinerate all waste
        \item Dump waste in oceans
        \item Store waste underground
    \end{parts}

\end{enumerate}

\end{document}
